\chapter{Sensibilidades del XVA por diferenciación adjunta}\label{cap:xva-aad}

% -------------------------------------------------------------------------------------------------
\section{Propósito}\label{sec:cap16-proposito}

El capítulo~\ref{cap:xva} calcula el CVA, el DVA y el FVA de una contraparte con una simulación de Monte Carlo. Una
mesa que quiera cubrir ese número necesita saber cuánto cambia si se mueve cada insumo: cada pilar de la curva
TIIE-F, de la curva SOFR y de la curva MXN-USD, el tipo de cambio, las volatilidades, la intensidad de cada tramo de
la curva de crédito de la contraparte y de la propia, y el diferencial de fondeo. En el snapshot sembrado son
48 insumos. Por choque y reprecio cada sensibilidad cuesta dos valuaciones completas; por \emph{diferenciación
adjunta} (AAD, \emph{adjoint algorithmic differentiation}) salen todas juntas por el costo de unas pocas
valuaciones, que es lo que este capítulo documenta.

El código calcula cuatro cosas. Primero, las derivadas de CVA, DVA y FVA (y, por combinación lineal, del total) respecto
de los 48 insumos, con la biblioteca XAD, sobre \emph{las mismas trayectorias, malla y semilla} que el XVA del
capítulo~\ref{cap:xva} (secciones~\ref{sec:cap16-insumos}--\ref{sec:cap16-cinta}). Segundo, el paso de las
intensidades de crédito a la sensibilidad a \emph{cada cotización de CDS} con el jacobiano del bootstrap
(sección~\ref{sec:cap16-credito}). Tercero, la derivada \emph{total} respecto de TIIE-F, SOFR y el spot, con la curva
MXN-USD recalibrada, por regla de la cadena con las parciales que la AAD ya entrega
(sección~\ref{sec:cap16-parcial}). Cuarto, la conversión de escalas que hace la ruta
\texttt{GET /api/xva/sensitivities} (sección~\ref{sec:cap16-implementacion}) y que alimenta la sección
«Sensibilidades» de la pantalla XVA. El cálculo es informativo, como todo el XVA: no entra al valor presente ni a
la contabilidad.

Lo que \emph{no} hace importa tanto como lo que hace, y se repite en la sección~\ref{sec:cap16-limites}: la AAD entrega
derivadas \textbf{parciales} (al mover un pilar no se recalibra ni la curva MXN-USD ni el crédito) y la ruta añade la
\textbf{total} con MXN-USD recalibrada, que es la que sirve para cubrir y puede tener el signo contrario
(sección~\ref{sec:cap16-parcial}); el crédito no se recalibra en ninguna. Son derivadas de primer orden, sin sensibilidad
a correlaciones ni al PFE y sin error estándar, y su costo (sección~\ref{sec:cap16-costo}) es de unas 7 a 9
valuaciones a igual número de hilos, de 15 a 23 veces el tiempo de una valuación en la ruta (la valuación usa todos
los hilos y la adjunta dos) y de decenas de segundos en conjuntos largos.

% -------------------------------------------------------------------------------------------------
\section{Notación}\label{sec:cap16-notacion}

Los símbolos comunes están en el apéndice~\ref{ap:simbolos}, y los del XVA ($t_j$, $V$, $\mathcal B$, $S_c$, $S_o$,
$\mathrm{LGD}$, $s_f$, $\mathrm{EPE}^*$, $\mathrm{ENE}^*$, $n$, $\mathcal G$) son los del capítulo~\ref{cap:xva}. Los
símbolos propios de este capítulo son:
\begin{itemize}
  \item $\boldsymbol\vartheta=(\vartheta_1,\dots,\vartheta_m)$: vector de \emph{insumos activos}, los que se derivan
        (tabla~\ref{tab:cap16-insumos}); $m=48$ en MX-DEMO. Se usa $\vartheta$ y no $\theta$ porque este último es la
        deriva $\hwtheta$ de Hull--White del capítulo~\ref{cap:hull-white}.
  \item $z^{\mathrm M}_k$, $z^{\mathrm U}_k$, $z^{\mathrm{col}}_k$: tasa cero (anual, composición continua) del pilar
        $k$ de las curvas TIIE-F, SOFR y MXN-USD (la curva de descuento con colateral \USD, subíndice \emph{col}).
  \item $h^{\mathrm c}_i$, $h^{\mathrm o}_i$: intensidad del tramo $i$ de la curva de crédito de la contraparte y de la
        propia; $v_i$ el vértice donde termina (años) y $J$ el número de tramos (en el capítulo~\ref{cap:xva} son
        $h_j$ y $v_j$). $\ell_i(t)$: tiempo vivido en el tramo $i$ hasta $t$.
  \item $\omega_p$: las normales de la trayectoria $p$ (cinco por día, de un \texttt{mt19937\_64} sembrado con
        (semilla, $p$)); \emph{no dependen de $\boldsymbol\vartheta$}. $V_{pj}$ y $\mathcal B_{pj}$: valor del
        conjunto y cuenta bancaria en la trayectoria $p$ y la fecha $t_j$ de la malla de integración $\mathcal G$.
  \item $w^{\mathrm c}_j$, $w^{\mathrm o}_j$, $w^{\mathrm f}_j$: pesos de la fecha $t_j$ en CVA, DVA y FVA
        (ecuación~\eqref{eq:cap16-pesos}); $\widehat{\mathrm{CVA}}_p$, $\widehat{\mathrm{DVA}}_p$,
        $\widehat{\mathrm{FVA}}_p$: aporte de la trayectoria $p$ (suman, divididos entre $n$, a las medidas).
  \item $\mathbf 1_{\{\cdot\}}$: indicador. $u_k$, $\bar u_k$: variable intermedia del programa y su adjunto
        ($\bar u_k=\partial\Phi/\partial u_k$ para la salida que se siembra). $\boldsymbol\kappa$: \emph{parte común}
        (variables que dependen de $\boldsymbol\vartheta$ y no de $\omega_p$).
  \item $s_k$, $J_{ki}=\partial h_i/\partial s_k$: cotización del CDS de plazo $T_k$ y jacobiano del bootstrap
        (adimensional).
  \item $\Delta_i$, $\widetilde D^{X}_i$: escala que lleva una derivada a la unidad del reporte y derivada reportada
        (ecuación~\eqref{eq:cap16-unidades}); $\DVO^{\mathrm{par}}_X$: DV01 paralelo de una curva.
  \item $C_{\mathrm{adj}}$, $C_{\mathrm{ef}}$: constante del costo de una pasada adjunta frente a la evaluación del
        programa y su valor medido aquí (la AAD frente a una valuación); $M_{\boldsymbol\kappa}$, $M_{\mathrm{tray}}$:
        memoria de la cinta de la parte común y de una trayectoria (bytes).
\end{itemize}

% -------------------------------------------------------------------------------------------------
\section{Modelo y fórmulas}\label{sec:cap16-modelo}

\begin{sloppypar}
\subsection{Qué se deriva y respecto de qué}\label{sec:cap16-insumos}

Se derivan el CVA, el DVA y el FVA del conjunto de compensación de una contraparte,
ecuaciones~(\ref{eq:cap16-pesos}) y (\ref{eq:cap16-contribuciones}), respecto de los insumos de la
tabla~\ref{tab:cap16-insumos}. El total es $\mathrm{CVA}-\mathrm{DVA}+\mathrm{FVA}$ y se deriva por linealidad:
\begin{equation}\label{eq:cap16-total-lineal}
  \frac{\partial\,\mathrm{Total}}{\partial\vartheta_i}=\frac{\partial\,\mathrm{CVA}}{\partial\vartheta_i}
  -\frac{\partial\,\mathrm{DVA}}{\partial\vartheta_i}+\frac{\partial\,\mathrm{FVA}}{\partial\vartheta_i}.
\end{equation}
La columna \texttt{dva} de la API es $\partial\mathrm{DVA}/\partial\vartheta_i$ (positiva cuando el DVA crece con el
insumo) y la \texttt{total} resta esa columna, como en la ecuación~\eqref{eq:cap16-total-lineal}.
\begin{table}[ht]
\centering\small
\begin{tabular}{@{}>{\raggedright\arraybackslash}p{3.5cm}>{\raggedright\arraybackslash}p{4.4cm}r>{\raggedright\arraybackslash}p{3.6cm}@{}}
\toprule
Grupo & Insumo $\vartheta_i$ & Cantidad & Unidad del reporte \\
\midrule
Curva TIIE-F & $z^{\mathrm M}_k$, un pilar por cotización & 10 & \MXN\ por \pb \\
Curva SOFR & $z^{\mathrm U}_k$ & 10 & \MXN\ por \pb \\
Curva MXN-USD (si el snapshot la trae) & $z^{\mathrm{col}}_k$ & 9 & \MXN\ por \pb \\
Tipo de cambio & $\Spot$ & 1 & \MXN\ por 1~\% del spot \\
Volatilidad de \USD/\MXN & $\sigma_X$ & 1 & \MXN\ por punto ($0.01$) \\
Hull--White & $\sigma_{\mathrm M}$, $\sigma_{\mathrm U}$ & 2 & \MXN\ por \pb\ de volatilidad normal \\
Hull--White & $a_{\mathrm M}$, $a_{\mathrm U}$ & 2 & \MXN\ por $0.01$ \\
Crédito de la contraparte & $h^{\mathrm c}_i$ (y cotizaciones $s_k$) & 6 & \MXN\ por \pb \\
Crédito propio & $h^{\mathrm o}_i$ (y cotizaciones $s_k$) & 6 & \MXN\ por \pb \\
Fondeo & $s_f$ & 1 & \MXN\ por \pb \\
\bottomrule
\end{tabular}
\caption{Insumos activos de \codigo{include/qf/xva/exposure_engine.hpp}{XvaGreeks} (cantidades de MX-DEMO, $m=48$) y
unidad en que la API los reporta.}\label{tab:cap16-insumos}
\end{table}

\emph{No son insumos}: las correlaciones $\rho_{\mathrm{MU}},\rho_{\mathrm{MX}},\rho_{\mathrm{UX}}$ (el modelo las usa solo
en la factorización de Cholesky, que no se deriva respecto de ellas), las LGD, los nocionales, las tasas fijas, las
fechas y las bases de días de las operaciones, las normales $\omega_p$, el número de trayectorias, la semilla y la
malla. Las tasas cero de la tabla~\ref{tab:cap16-insumos} son las tasas del \emph{pilar} (una por cotización de
depósito u OIS del snapshot, o por punto forward y base en MXN-USD), no las cotizaciones mismas
(sección~\ref{sec:cap16-parcial}).

El reporte multiplica cada derivada por la escala $\Delta_i$ que la lleva a su unidad,
\begin{equation}\label{eq:cap16-unidades}
  \widetilde D^{X}_i=\Delta_i\,\frac{\partial X}{\partial\vartheta_i},\qquad
  \Delta_i=\begin{cases}
    10^{-4} & \text{tasas cero, }\sigma_{\mathrm M},\sigma_{\mathrm U},\ h^{\mathrm c}_i,\ h^{\mathrm o}_i,\ s_f\ (1~\pb),\\
    0.01\,\Spot & \text{spot (1~\% del spot, en \MXN\ por \USD)},\\
    0.01 & \sigma_X\ (\text{un punto de volatilidad}),\ a_{\mathrm M},\ a_{\mathrm U},
  \end{cases}
\end{equation}
con $X\in\{\mathrm{CVA},\mathrm{DVA},\mathrm{FVA}\}$ y $\widetilde D^{\mathrm{Total}}=\widetilde D^{\mathrm{CVA}}-\widetilde D^{\mathrm{DVA}}+\widetilde D^{\mathrm{FVA}}$. El DV01 paralelo de una curva es
la suma de sus pilares, $\DVO^{\mathrm{par}}_X=\sum_k\widetilde D^{X}_{z_k}$, porque con interpolación lineal y
extrapolación plana desplazar todos los pilares la misma cantidad desplaza toda la curva esa cantidad. Las
conversiones están en \codigo{desk/api/risk/xva_sens.py}{_row}, \codigo{desk/api/risk/xva_sens.py}{_curves} y
\codigo{desk/api/risk/xva_sens.py}{VOLS}; los rótulos de unidad los pone la interfaz
(\codigo{desk/ui/src/screens/xva/XvaSensitivitiesPanel.tsx}{UNIT_VOLS}) y la respuesta no los repite.

\subsection{Modo adjunto}\label{sec:cap16-adjunto}

Sea $\Phi:\mathbb R^m\to\mathbb R^q$ el programa que calcula las salidas, descompuesto en operaciones elementales
$u_k=\varphi_k(u_j:\ j\prec k)$ con $u_i=\vartheta_i$ para $i\le m$ y salidas $\Phi_l=u_{L_l}$. El modo
\emph{directo} propaga una dirección de entrada por pasada: un insumo, una pasada. El modo \emph{adjunto} recorre el
programa grabado hacia atrás y entrega la derivada de una salida respecto de \emph{todos} los insumos en una sola
pasada \citep{griewank2008,naumann2011}:
\begin{equation}\label{eq:cap16-adjunto}
  \bar u_{L_l}=1,\qquad
  \bar u_j=\sum_{k:\ j\prec k}\bar u_k\,\frac{\partial\varphi_k}{\partial u_j},\qquad
  \frac{\partial\Phi_l}{\partial\vartheta_i}=\bar u_i\quad(i\le m).
\end{equation}
Una pasada hacia atrás cuesta una constante $C_{\mathrm{adj}}$ pequeña, independiente de $m$, por la evaluación de $\Phi$
\citep{griewank2008}, mientras que la diferencia central cuesta $2m$ evaluaciones:
\begin{equation}\label{eq:cap16-costo}
  \mathrm{coste}(\nabla\Phi_l)\ \le\ C_{\mathrm{adj}}\,\mathrm{coste}(\Phi)\qquad\text{frente a}\qquad
  \mathrm{coste}(\text{choque central})=2m\,\mathrm{coste}(\Phi).
\end{equation}
Es la razón por la que el modo adjunto domina cuando hay pocas salidas y muchos insumos
\citep{giles2005,capriotti2011}, que es exactamente el caso del XVA, y por la que se ha aplicado a la
gestión del riesgo de crédito de contraparte \citep{capriotti2011b} (\citealp{capriotti2024} revisa los quince años de
uso de la AAD en finanzas). Aquí hay $q=3$ salidas (CVA, DVA y FVA) y la biblioteca (XAD) las propaga a la vez: cada variable
activa lleva un \emph{vector} de adjuntos de tamaño 4 (\texttt{xad::AReal<double, 4>}, el alias \texttt{AD} de
\codigo{src/xva/exposure_engine.cpp}{ExposureEngine::xvaSensitivities}; la biblioteca solo trae instancias de
tamaño 1, 2 y 4, así que la cuarta componente queda libre), y una pasada hacia atrás con las semillas $(1,0,0)$, $(0,1,0)$ y
$(0,0,1)$ en las tres primeras componentes entrega los tres gradientes. Con un adjunto de tamaño 1 (solo el total) la
cinta es un 20~\% más rápida, medido en la sonda del diseño, pero no separa CVA, DVA y FVA y no se usó.

La razón de costo medida del programa completo es $C_{\mathrm{ef}}\approx9.3$ con un solo hilo
(sección~\ref{sec:cap16-costo}); no es el $C_{\mathrm{adj}}$ teórico porque incluye grabar la cinta y el tráfico de memoria. Con
ella, el modo adjunto gana a los choques cuando $2m>9.3$, es decir, desde 5 insumos \emph{a igual número de hilos};
con los 48 insumos del snapshot la ventaja es de $96/9.3\approx10$ veces.

\subsection{El estimador trayectoria a trayectoria}\label{sec:cap16-estimador}

Con $\omega_p$ fijas, el CVA, el DVA y el FVA del capítulo~\ref{cap:xva} son promedios de funciones de $\boldsymbol\vartheta$.
Sean, para $j\ge1$,
\begin{equation}\label{eq:cap16-pesos}
\begin{gathered}
  w^{\mathrm c}_j=\mathrm{LGD}_c\bigl[S_c(t_{j-1})-S_c(t_j)\bigr],\qquad
  w^{\mathrm o}_j=\mathrm{LGD}_o\bigl[S_o(t_{j-1})-S_o(t_j)\bigr],\\[2pt]
  w^{\mathrm f}_j=s_f\,S_c(t_j)S_o(t_j)\,(t_j-t_{j-1}),
\end{gathered}
\end{equation}
y, como $V^+-V^-=V$, los aportes de la trayectoria $p$ (la ecuación del error estándar del capítulo~\ref{cap:xva}, para
el CVA, y sus análogas)
\begin{equation}\label{eq:cap16-contribuciones}
  \widehat{\mathrm{CVA}}_p=\sum_{j\ge1}w^{\mathrm c}_j\,\frac{V^+_{pj}}{\mathcal B_{pj}},\qquad
  \widehat{\mathrm{DVA}}_p=\sum_{j\ge1}w^{\mathrm o}_j\,\frac{V^-_{pj}}{\mathcal B_{pj}},\qquad
  \widehat{\mathrm{FVA}}_p=\sum_{j\ge1}w^{\mathrm f}_j\,\frac{V_{pj}}{\mathcal B_{pj}},
\end{equation}
de modo que $X=\tfrac1n\sum_p\widehat X_p$ para cada medida $X$. (La suma sobre fechas queda \emph{dentro} de la
trayectoria; el capítulo~\ref{cap:xva} promedia primero sobre trayectorias para obtener $\mathrm{EPE}^*$ y
$\mathrm{ENE}^*$. Son la misma cantidad con otro orden de sumas.) Como $\mathcal B>0$, $V^+/\mathcal B=(V/\mathcal B)^+$,
y el estimador de la derivada es el promedio de las derivadas de cada aporte:
\begin{equation}\label{eq:cap16-pathwise}
  \frac{\partial\widehat{\mathrm{CVA}}_p}{\partial\vartheta_i}=\sum_{j\ge1}\Bigl[
  w^{\mathrm c}_j\,\mathbf 1_{\{V_{pj}>0\}}\,\frac{\partial}{\partial\vartheta_i}\Bigl(\frac{V_{pj}}{\mathcal B_{pj}}\Bigr)
  +\frac{\partial w^{\mathrm c}_j}{\partial\vartheta_i}\,\frac{V^+_{pj}}{\mathcal B_{pj}}\Bigr],\qquad
  \widehat{\partial_i\mathrm{CVA}}=\frac1n\sum_{p=1}^{n}\frac{\partial\widehat{\mathrm{CVA}}_p}{\partial\vartheta_i}.
\end{equation}
Para el DVA el indicador es $-\mathbf 1_{\{V_{pj}<0\}}$ (porque $V^-=\max(-V,0)$) y para el FVA no hay indicador
(la función es diferenciable) y la regla del producto actúa sobre los tres factores de $w^{\mathrm f}_j\,V_{pj}/\mathcal B_{pj}$. Los dos sumandos
de la ecuación~\eqref{eq:cap16-pathwise} se reparten los insumos: $V_{pj}$ y $\mathcal B_{pj}$ dependen de las curvas
de tasa, el spot y los parámetros del modelo, y \emph{no} del crédito ni del fondeo, de modo que para los primeros el segundo sumando es cero, y
para las intensidades $h^{\mathrm c}_i$, $h^{\mathrm o}_i$ y el diferencial $s_f$ (que solo están en los pesos) es cero el
primero. Es lo que hace \codigo{src/xva/exposure_engine.cpp}{ExposureEngine::xvaSensitivities}: el indicador es la derivada de \texttt{max(V, 0.0)} de XAD, que vale $1$ si
$V>0$, $0$ si $V<0$ y $\tfrac12$ en el empate (convención de XAD 2.2.0).

\paragraph{Por qué es insesgado.} El estimador es el \emph{método de la derivada por trayectoria}
\citep{broadie1996,glasserman2003}: se deriva el aporte de cada trayectoria con las normales fijas y se promedia.
Intercambiar derivada y esperanza, $\partial_{\vartheta}\E[\widehat X_p]=\E[\partial_{\vartheta}\widehat X_p]$, es válido si
(i) para cada $\omega_p$ la función $\vartheta\mapsto\widehat X_p(\vartheta)$ es continua y, con probabilidad uno,
diferenciable en el punto, y (ii) es lipschitziana con una constante aleatoria de esperanza finita
(convergencia dominada). (i) se cumple porque $\max(\cdot,0)$ solo no es diferenciable donde $V_{pj}=0$, un evento de
probabilidad cero en las fechas donde $V_{pj}$ tiene densidad; en las fechas en que $V\equiv0$ (después de pagado el
último flujo) el valor no depende de $\boldsymbol\vartheta$ y su derivada es cero, cualquiera que sea la convención del
empate. $\max(\cdot,0)$ es 1-lipschitziana y $V/\mathcal B$ es una función suave de $\boldsymbol\vartheta$ para cada $\omega_p$
(exponenciales de estados gaussianos), por lo que (ii) es la condición estándar del método. El código \emph{no
verifica} (ii); lo que sí hace el repositorio es comprobar el resultado (sección~\ref{sec:cap16-ej-vega}: la media de
lotes independientes cae sobre la derivada exacta de una forma cerrada). Además, con las normales fijas el
estimador es exactamente la derivada del \emph{estimador de Monte Carlo} $X(\boldsymbol\vartheta)=\tfrac1n\sum_p\widehat X_p$,
y por eso coincide, en el límite del paso, con la diferencia central que reutiliza las mismas trayectorias; esa
coincidencia es la prueba de corrección de las secciones~\ref{sec:cap16-ej-choque} y~\ref{sec:cap16-validacion}.

\paragraph{Por qué el PFE no se deriva.} El $\mathrm{PFE}_\alpha(t_i)$ es el elemento de rango $\lceil\alpha n\rceil$
de los $V^+_p(t_i)$ ordenados: no es un promedio de una función de la trayectoria, y su derivada trayectoria a
trayectoria es la de la \emph{única} trayectoria que ocupa ese rango, no un promedio. No cumple la estructura de la
ecuación~\eqref{eq:cap16-pathwise} y exigiría otro estimador (de razón de verosimilitud o suavizado); este corte no
lo implementa y la API no publica sensibilidades del PFE.

\subsection{La cinta por trayectoria}\label{sec:cap16-cinta}

Grabar la simulación completa en una sola cinta es inviable: cada trayectoria añade unos 43 MB de cinta
(sección~\ref{sec:cap16-costo}), así que 2000 trayectorias serían del orden de 86 GB. El código aprovecha que el
programa tiene dos partes. La \emph{parte común} $\boldsymbol\kappa(\boldsymbol\vartheta)$ reúne todo lo que depende
de los insumos y no de las normales: las tablas por día de la malla (descuentos, $y_c$, $m_c$, $M_c$, factores
$e^{-a_c d/365}$, $q$), la factorización de Cholesky de la covarianza del paso, el spot, las tablas del contexto de
las operaciones y los pesos de la ecuación~\eqref{eq:cap16-pesos}. La \emph{parte de trayectoria} calcula
$(\widehat{\mathrm{CVA}}_p,\widehat{\mathrm{DVA}}_p,\widehat{\mathrm{FVA}}_p)=G(\boldsymbol\kappa,\omega_p)$. Por la regla de la cadena,
\begin{equation}\label{eq:cap16-parte-comun}
  \frac{\partial X}{\partial\vartheta_i}=\sum_r\frac{\partial X}{\partial\kappa_r}\,\frac{\partial\kappa_r}{\partial\vartheta_i},\qquad
  \frac{\partial X}{\partial\kappa_r}=\frac1n\sum_{p=1}^{n}\frac{\partial\widehat X_p}{\partial\kappa_r},
\end{equation}
y el patrón (en \codigo{src/xva/exposure_engine.cpp}{forEachPathChunk} y \codigo{src/xva/exposure_engine.cpp}{ExposureEngine::xvaSensitivities}) consiste en grabar $\boldsymbol\kappa$ una vez y, para cada trayectoria, grabar solo su parte, barrerla hacia
atrás hasta $\boldsymbol\kappa$ (los adjuntos de $\boldsymbol\kappa$ se \emph{acumulan} en sus ranuras) y borrarla:
\begin{enumerate}
  \item registrar los insumos y empezar a grabar; calcular $\boldsymbol\kappa$ (tablas y pesos) y guardar la posición $\pi$ de la cinta;
  \item para cada trayectoria $p$ del hilo: grabar $G(\boldsymbol\kappa,\omega_p)$ (\texttt{simulatePath}, \texttt{computeBonds},
        \texttt{setValue} en cada fecha de la malla y la acumulación de las tres salidas), registrar las tres salidas y
        sembrar $(1,0,0)$, $(0,1,0)$ y $(0,0,1)$; barrer hacia atrás \emph{solo hasta $\pi$}
        (\texttt{computeAdjointsTo}); poner a cero las derivadas posteriores a $\pi$ (\texttt{clearDerivativesAfter}); y
        regresar la cinta a $\pi$ (\texttt{resetTo});
  \item un último barrido de $\boldsymbol\kappa$ a los insumos (\texttt{computeAdjoints}) entrega $\partial X/\partial\vartheta_i$.
\end{enumerate}
La cinta de un hilo no crece con el número de trayectorias:
\begin{equation}\label{eq:cap16-memoria}
  M_{\mathrm{cinta}}=M_{\boldsymbol\kappa}+M_{\mathrm{tray}}\quad(\text{independiente de }n),\qquad\text{frente a}\qquad
  M_{\boldsymbol\kappa}+n\,M_{\mathrm{tray}}.
\end{equation}
Ninguna variable activa de una trayectoria puede sobrevivir a \texttt{resetTo}; por eso los datos de cada
trayectoria se declaran dentro del bucle. Con hilos, cada uno tiene su propia cinta, su copia de los insumos y su
propia parte común, y guarda sus adjuntos parciales indexados por su primera trayectoria; se suman en ese orden, así
que el resultado no depende de cuándo termine cada hilo y el número de hilos solo cambia el orden de redondeo de las
sumas. La API usa como máximo 2 hilos (sección~\ref{sec:cap16-costo}). Derivar por trayectoria y recorrer hacia atrás cada una es la técnica de
\citet{giles2005} y \citet{capriotti2011}; repartirla en hilos con una cinta por hilo, la de la AAD sobre simulaciones
paralelas \citep{savine2018}.

\subsection{Crédito: de intensidades a cotizaciones de CDS}\label{sec:cap16-credito}

Las probabilidades de supervivencia entran al motor como función de las intensidades por tramo, que son insumos
activos. Con vértices $0=v_0<v_1<\dots<v_J$ y extrapolación plana de la última intensidad, la plantilla
\codigo{include/qf/xva/piecewise_hazard_curve.hpp}{survivalFromHazards} calcula (es la
ecuación de supervivencia del capítulo~\ref{cap:xva} con el tiempo vivido en cada tramo $\ell_i$ escrito explícito)
\begin{equation}\label{eq:cap16-supervivencia-activa}
  S(t)=\exp\Bigl(-\sum_{i=1}^{J}h_i\,\ell_i(t)\Bigr),\quad
  \ell_i(t)=\begin{cases}[\min(t,v_i)-v_{i-1}]^+ & i<J,\\ [\,t-v_{J-1}]^+ & i=J,\end{cases}\quad
  \frac{\partial S(t)}{\partial h_i}=-\ell_i(t)\,S(t).
\end{equation}
Como el perfil de exposición descontado no depende de las intensidades, las derivadas respecto de ellas son exactas
dado el perfil (no llevan derivada de Monte Carlo):
\begin{equation}\label{eq:cap16-dcva-dh}
\begin{aligned}
  \frac{\partial\,\mathrm{CVA}}{\partial h^{\mathrm c}_i}&=-\mathrm{LGD}_c\sum_{j\ge1}\mathrm{EPE}^*(t_j)\bigl[\ell_i(t_{j-1})S_c(t_{j-1})-\ell_i(t_j)S_c(t_j)\bigr],\\
  \frac{\partial\,\mathrm{FVA}}{\partial h^{\mathrm c}_i}&=-s_f\sum_{j\ge1}\bigl[\mathrm{EPE}^*(t_j)-\mathrm{ENE}^*(t_j)\bigr]\,\ell_i(t_j)\,S_c(t_j)S_o(t_j)\,(t_j-t_{j-1}),\qquad
  \frac{\partial\,\mathrm{FVA}}{\partial s_f}=\frac{\mathrm{FVA}}{s_f},
\end{aligned}
\end{equation}
y el DVA y el FVA respecto de $h^{\mathrm o}_i$ son análogos con $\mathrm{ENE}^*$, $S_o$ y $\mathrm{LGD}_o$. El código
no usa estas fórmulas: obtiene los mismos números de la cinta, y la sección~\ref{sec:cap16-ej-credito} las usa
como comprobación independiente. Los insumos «de crédito» que ve la mesa, sin embargo, son las cotizaciones de CDS
$s_k$, no las intensidades. Cada intensidad es una función de las cotizaciones por el bootstrap de la
ecuación de CDS del capítulo~\ref{cap:xva}: $G_k(h_1,\dots,h_k;s_k)=V_{\mathrm{CDS}}(T_k,s_k;h_1,\dots,h_k)=0$,
$k=1,\dots,J$. Derivando respecto de $s_{k'}$ y recordando que $\partial V_{\mathrm{CDS}}/\partial s_k=-\Nom\,\Ann_{\mathrm{cds},k}$,
\begin{equation}\label{eq:cap16-jacobiano}
  \sum_{i\le k}\frac{\partial G_k}{\partial h_i}\,J_{k'i}=\delta_{kk'}\,\Nom\,\Ann_{\mathrm{cds},k},\qquad
  J_{k'i}=\frac{\partial h_i}{\partial s_{k'}}=0\ \ \text{si }i<k',
\end{equation}
un sistema triangular: $h_i$ depende de las cotizaciones de plazo menor o igual, y mover la cotización $k'$
mueve la intensidad del tramo $k'$ y, a través de las supervivencias que cada CDS posterior hereda, las de todos los
tramos posteriores. Por la regla de la cadena,
\begin{equation}\label{eq:cap16-cadena}
  \frac{\partial X}{\partial s_k}=\sum_{i=k}^{J}\frac{\partial X}{\partial h_i}\,J_{ki}.
\end{equation}
\textbf{El jacobiano no se obtiene resolviendo la ecuación~\eqref{eq:cap16-jacobiano} ni por AAD del bootstrap}
(Brent no tiene derivada útil): \codigo{desk/api/risk/xva_sens.py}{_credit} \emph{reprecia} el bootstrap con
\codigo{desk/api/risk/credit_curves.py}{build_curve}, moviendo la cotización $k$ en $\pm\varepsilon$ con $\varepsilon=10^{-6}$ ($0.01$~\pb;
\codigo{desk/api/risk/xva_sens.py}{CDS_BUMP}) y restando,
\begin{equation}\label{eq:cap16-jacobiano-reprecio}
  J_{ki}\approx\frac{h_i(s_k+\varepsilon)-h_i(s_k-\varepsilon)}{2\varepsilon},
\end{equation}
que cuesta $2J$ bootstraps por entidad. La respuesta de la ruta declara el método en el campo \texttt{cds\_jacobian}.
Con una curva plana o supuesta (sin cotizaciones) no hay cotizaciones que derivar: la ruta devuelve un solo tramo, con
\texttt{to\_years} nulo, y la derivada es respecto de esa intensidad.

\subsection{Derivadas parciales y derivada total con MXN-USD recalibrada}\label{sec:cap16-parcial}

Los insumos $\boldsymbol\vartheta$ son los del \emph{motor}; los datos de mercado son otros. La curva MXN-USD se calibra
(\codigo{desk/api/market.py}{_calibrate_csa}) con \texttt{bootstrap\_xccy\_discount} a partir de puntos forward, cotizaciones de base, el spot y las curvas
SOFR y TIIE-F; las intensidades salen de las cotizaciones de CDS descontadas con la curva TIIE-F; y las curvas de tasa
son ceros de pilares que a su vez salen de cotizaciones de depósito y OIS. La derivada que entrega la AAD es la parcial
$\partial X/\partial\vartheta_i$ \emph{sin recalibrar nada}: al mover un pilar de TIIE-F, las tasas cero de MXN-USD
y las intensidades de crédito se quedan donde estaban. La derivada \emph{total} respecto de ese pilar sumaría los
efectos de la recalibración,
\begin{equation}\label{eq:cap16-total}
  \frac{dX}{dz^{\mathrm M}_k}=\frac{\partial X}{\partial z^{\mathrm M}_k}
  +\sum_{l}\frac{\partial X}{\partial z^{\mathrm{col}}_l}\,\frac{\partial z^{\mathrm{col}}_l}{\partial z^{\mathrm M}_k}
  +\sum_{i}\frac{\partial X}{\partial h^{\mathrm c}_i}\,\frac{\partial h^{\mathrm c}_i}{\partial z^{\mathrm M}_k}
  +\sum_{i}\frac{\partial X}{\partial h^{\mathrm o}_i}\,\frac{\partial h^{\mathrm o}_i}{\partial z^{\mathrm M}_k},
\end{equation}
donde el primer sumando es lo que entrega la AAD, el segundo (la recalibración de MXN-USD) usa las parciales respecto de
los pilares de MXN-USD que la AAD \emph{también} entrega, y los dos últimos (el crédito) exigen el jacobiano del
bootstrap de CDS respecto de la curva de descuento.

La ruta calcula la total con los dos primeros sumandos, \emph{sin simulación adicional}, para cada pilar de TIIE-F y de
SOFR y para el spot (\codigo{desk/api/risk/xva_sens.py}{_total}):
\begin{equation}\label{eq:cap16-total-api}
  \frac{dX}{d\vartheta_i}\bigg|_{\mathrm{MXN\text{-}USD}}=\frac{\partial X}{\partial\vartheta_i}
  +\sum_{l}\frac{\partial X}{\partial z^{\mathrm{col}}_l}\,\frac{\partial z^{\mathrm{col}}_l}{\partial\vartheta_i},\qquad
  \frac{\partial z^{\mathrm{col}}_l}{\partial\vartheta_i}\approx
  \frac{z^{\mathrm{col}}_l(\vartheta_i+\epsilon_i)-z^{\mathrm{col}}_l(\vartheta_i-\epsilon_i)}{2\epsilon_i},
\end{equation}
con $\vartheta_i\in\{z^{\mathrm M}_k,z^{\mathrm U}_k,\Spot\}$, $\epsilon_i$ igual a 1~\pb\ de tasa o a $0.1\,\%$ del spot,
y $z^{\mathrm{col}}_l(\cdot)$ la tasa del pilar $l$ de MXN-USD que resulta de recalibrarla con sus puntos forward y su
base \emph{fijos} sobre las curvas y el spot chocados: es la función que el desk ya usa para escenarios
(\codigo{desk/api/market.py}{MarketContext.shifted}), no una calibración nueva. Las derivadas
$\partial X/\partial\vartheta_i$ y $\partial X/\partial z^{\mathrm{col}}_l$ son las de la AAD, sin escalar; el reporte las
lleva a su unidad con la ecuación~\eqref{eq:cap16-unidades}. El jacobiano cuesta $2\times(10+10+1)=42$ recalibraciones
(unos 1.6~ms cada una, 0.07~s en total frente a los 5--12 s de la AAD) y ninguna simulación. Cada fila de TIIE-F, de SOFR
y del spot trae entonces las dos derivadas: la parcial (\texttt{cva}, \texttt{dva}, \texttt{fva}, \texttt{total}, con
MXN-USD fija) y la total (\texttt{recalibrated}, con MXN-USD recalibrada), y el DV01 paralelo de cada curva las trae
también, como suma de sus pilares. MXN-USD no tiene total: sus tasas cero son lo que se recalibra, así que su riesgo de
mercado ya está en la total de TIIE-F, SOFR y el spot (y en las cotizaciones de puntos forward y de base, que se
mantienen fijas y no se derivan); para el resto de los insumos (volatilidades, crédito, fondeo) la calibración de
MXN-USD no interviene y la total es la parcial, y sin MXN-USD las dos coinciden. Los dos últimos sumandos de la
ecuación~\eqref{eq:cap16-total} \emph{no} se calculan: no se recalibra el bootstrap de CDS al mover TIIE-F. La
sección~\ref{sec:cap16-ej-parcial} mide qué tan distinta es la total de la parcial en un cross currency swap: tiene
incluso el signo contrario. Para cubrir con instrumentos de TIIE-F, de SOFR o del tipo de cambio es la total la que
dice cuánto cambia el XVA cuando esa tasa se mueve en el mercado; la parcial sirve para ver de dónde viene el riesgo.

\begin{diferencia}
  Una sensibilidad de mercado completa derivaría la cadena
  $\text{cotizaciones}\to\text{curvas}\to\text{parámetros del modelo}\to\text{XVA}$ entera, incluida la calibración (la
  AAD lo permite, por ejemplo con el teorema de la función implícita en cada paso de calibración). Aquí se deriva el motor (que recibe pilares, parámetros e intensidades ya
  calibrados) y se declara qué falta: (i) la AAD entrega la derivada \emph{parcial} y la ruta añade la total con MXN-USD recalibrada, ecuación~\eqref{eq:cap16-total-api},
pero no la del crédito (los dos últimos sumandos de la ecuación~\eqref{eq:cap16-total}); (ii) las curvas de tasa se
  derivan por tasa cero del pilar, no por cotización de depósito u OIS (para el crédito sí se pasa a cotizaciones, con
  el jacobiano del bootstrap por reprecio y no por AAD del bootstrap); (iii) $a_c$, $\sigma_c$ y $\sigma_X$ son
  insumos, no resultados de una calibración, y no se deriva respecto de la superficie de volatilidad de la que
  sale $\sigma_X$; (iv) toda la fórmula afín de los bonos y de las patas se graba en la cinta con el tipo genérico de
  XAD en vez de escribir a mano su adjunto (es el $70\,\%$ de la cinta; sección~\ref{sec:cap16-costo}); y (v) la
  derivada es la de la estimación discreta del capítulo~\ref{cap:xva} (malla de integración $\mathcal G$, extremo
  derecho), no la del XVA en tiempo continuo.
\end{diferencia}
\end{sloppypar}

% -------------------------------------------------------------------------------------------------
\section{Implementación}\label{sec:cap16-implementacion}

\begin{sloppypar}
\subsection{Un motor, dos escalares}\label{sec:cap16-escalar}

Las funciones del camino de exposición de \codigo{src/xva/exposure_engine.cpp}{ExposureEngine::simulate} (tablas por
día, \codigo{src/xva/exposure_engine.cpp}{buildGrid}, \codigo{src/xva/exposure_engine.cpp}{simulatePath},
\codigo{src/xva/exposure_engine.cpp}{computeBonds}, \codigo{src/xva/exposure_engine.cpp}{legValue} y
\codigo{src/xva/exposure_engine.cpp}{setValue}) son plantillas sobre el escalar \texttt{Real}: la misma fuente sirve
para \texttt{double} y para el tipo activo de XAD. No se duplicó la lógica y \emph{ningún número en \texttt{double} cambió}:
la prueba \codigo{tests/test_ExposureEngine.cpp}{GoldenValuesUnchangedByTemplating} (escrita antes de templar el motor)
compara con \texttt{EXPECT\_EQ} el CVA, DVA, FVA, perfiles y valores iniciales de siete operaciones con literales de 17 dígitos. Se
quedan en \texttt{double} lo que no es insumo: las normales, las fechas y días, las patas compiladas
(\codigo{src/xva/exposure_engine.cpp}{compileSet}), la validación y \codigo{src/xva/exposure_engine.cpp}{computeXva}.

Las curvas de entrada de la AAD son los pilares como insumos activos:
\codigo{src/xva/exposure_engine.cpp}{PillarCurve} guarda los plazos (\texttt{double}) y las tasas cero (\texttt{Real}), y su
descuento replica la rama \texttt{Linear} de la interpolación del capítulo~\ref{cap:curvas}: tasa constante fuera del
rango, interpolación lineal de las tasas cero en los pilares que encierran el plazo y $e^{-rT}$. En \texttt{double} se
sigue llamando a \texttt{YieldCurve} (cualquier interpolación); la AAD exige \texttt{Linear} y
\codigo{src/xva/exposure_engine.cpp}{requireLinear} lanza \texttt{std::invalid\_argument} si no es así (las curvas del desk lo
son). Son \textbf{tres excepciones} declaradas en comentarios del código a la regla de no duplicar:
la interpolación de unas 15 líneas, el bucle por trayectoria de la AAD (paralelo al de \texttt{simulate}, cuya
ejecución en \texttt{double} no debe cambiar) y la suma ponderada de las fechas, que repite la de \texttt{computeXva} en
otro orden (reutilizar \texttt{computeXva} habría movido sus números por redondeo).

\subsection{\texttt{xvaSensitivities}}\label{sec:cap16-xvasens}

\codigo{src/xva/exposure_engine.cpp}{ExposureEngine::xvaSensitivities} (declarada en
\codigo{include/qf/xva/exposure_engine.hpp}{ExposureEngine}, resultado en
\codigo{include/qf/xva/exposure_engine.hpp}{XvaSensitivities} y \codigo{include/qf/xva/exposure_engine.hpp}{XvaGreeks})
recibe el conjunto, las dos curvas de crédito con sus LGD, el diferencial de fondeo y el tope de hilos
(\texttt{maxThreads}, por omisión 2). Paso a paso:
\begin{enumerate}
  \item Valida las LGD, el diferencial y que las tres curvas de tasa sean \texttt{Linear}. Convierte las curvas de
        crédito a tramos con \codigo{src/xva/exposure_engine.cpp}{hazardsOf}: \texttt{PiecewiseHazardCurve} da sus
        vértices e intensidades y \texttt{FlatHazardRate}($\lambda$) un solo tramo, con vértice
        $\max(\text{último día}/365,\,1)$ (exactamente $e^{-\lambda t}$); cualquier otro tipo de curva lanza
        \texttt{std::invalid\_argument}.
  \item Arma el vector de insumos en este orden (que es el de \texttt{XvaGreeks}): pilares de TIIE-F, SOFR y
        MXN-USD (ninguno si el snapshot no trae la curva), spot, $(\sigma_X,\sigma_{\mathrm M},a_{\mathrm M},\sigma_{\mathrm U},a_{\mathrm U})$,
        intensidades de la contraparte y propias, y $s_f$.
  \item \codigo{src/xva/exposure_engine.cpp}{forEachPathChunk} reparte las trayectorias en trozos contiguos, uno por
        hilo (hasta \texttt{maxThreads}, por omisión 2). Cada hilo, con su cinta
        (\texttt{xad::Tape<double, 4>}), registra su copia de los insumos, graba $\boldsymbol\kappa$ y los pesos
        $w^{\mathrm c}_j,w^{\mathrm o}_j,w^{\mathrm f}_j$ (que ya llevan el $1/n$ del promedio) y ejecuta el patrón de
        la sección~\ref{sec:cap16-cinta}.
  \item Suma los parciales en el orden de las primeras trayectorias de cada trozo, y arma el valor:
        \texttt{value} es la suma de los aportes por trayectoria (CVA, DVA, FVA, total), con errores estándar
        calculados como en \texttt{computeXva}; coincide con \texttt{computeXva} hasta redondeo
        (diferencia relativa $\le10^{-12}$, probada).
\end{enumerate}
\codigo{src/xva/exposure_engine.cpp}{aadTapeBytes} y \codigo{src/xva/exposure_engine.cpp}{aadThreads} son
diagnósticos (pico de bytes de una cinta y número de cintas de la última llamada); no son API. Con
\texttt{QF\_WITH\_XAD=OFF} el método existe y lanza \texttt{std::logic\_error}.

\subsection{Binding, módulo de Python y ruta}\label{sec:cap16-ruta}

El módulo \texttt{qfxva} (\codigo{src/python_bindings/xva_bindings.cpp}{xva_sensitivities}) expone
\texttt{ExposureEngine.xva\_sensitivities(\dots, max\_threads=2)} (suelta el GIL), las clases \texttt{XvaGreeks} y
\texttt{XvaSensitivities} y la bandera \codigo{src/python_bindings/xva_bindings.cpp}{HAS_XAD}. La lógica de la ruta
está en \codigo{desk/api/risk/xva_sens.py}{compute}, que \emph{no contiene fórmulas financieras}:
\begin{enumerate}
  \item arma los mismos insumos que \texttt{GET /api/xva} reutilizando sus funciones (\texttt{build\_set},
        \texttt{model\_from\_snapshot}, \texttt{integration\_days}, parámetros y semilla) y obtiene el XVA de referencia
        llamando a la misma \texttt{xva.\_counterparty}; por eso el XVA que devuelve es \emph{igual a igual} al de
        \texttt{GET /api/xva} (el tiempo de esa llamada es \texttt{valuation\_elapsed\_s});
  \item llama a \texttt{xva\_sensitivities} sobre la malla de integración $\mathcal G$ (\texttt{elapsed\_s}) y
        convierte las unidades con la ecuación~\eqref{eq:cap16-unidades};
  \item compone las intensidades con el jacobiano de la ecuación~\eqref{eq:cap16-jacobiano-reprecio}
        (\codigo{desk/api/risk/xva_sens.py}{_credit}) y las parciales con el jacobiano de la calibración de MXN-USD,
        ecuación~\eqref{eq:cap16-total-api} (\codigo{desk/api/risk/xva_sens.py}{_total});
  \item cuenta los insumos de verdad (\texttt{n\_inputs}: las derivadas que entrega \texttt{qfxva}, 48 en MX-DEMO) y
        estima el costo por choques como $2\times n_{\mathrm{insumos}}\times\texttt{valuation\_elapsed\_s}$
        (\texttt{bump\_cost\_estimate\_s}: es una estimación, no una medición, y supone los choques con la valuación en
        todos los hilos), y añade a \texttt{assumptions} las notas \codigo{desk/api/risk/xva_sens.py}{PARTIAL_NOTE}
        (qué es la parcial y qué la total, y que el crédito no se recalibra en ninguna) y
        \codigo{desk/api/risk/xva_sens.py}{COST_NOTE}.
\end{enumerate}
La ruta es \codigo{desk/api/routes/xva.py}{xva_sensitivities_get}:
\texttt{GET /api/xva/sensitivities?snapshot=<id>\&counterparty=<nombre>}. Falta un parámetro $\to$ 400 con \texttt{field};
snapshot o contraparte inexistentes $\to$ 404; \texttt{HAS\_XAD} falso $\to$ 501; una contraparte sin swaps de tasa ni
cross currency swaps vivos $\to$ 200 con listas vacías y \texttt{empty\_reason}. La respuesta trae el DV01 paralelo por
curva, las filas por pilar, spot, volatilidades, CDS por plazo y fondeo, cada una con \texttt{cva}, \texttt{dva},
\texttt{fva} y \texttt{total} (la parcial) en \MXN\ por la unidad de la tabla~\ref{tab:cap16-insumos}, y los pilares de
TIIE-F y de SOFR, el spot y los DV01 paralelos de esas dos curvas con la misma fila en \texttt{recalibrated} (la total).
Con cotizaciones de CDS la respuesta trae la sensibilidad a cada cotización y no repite la intensidad por tramo
(\texttt{hazards} queda vacía); sin curva de CDS trae solo la del tramo único. La respuesta no lleva las unidades: las
pone la interfaz.

\subsection{Compilación, licencia y versión de XAD}\label{sec:cap16-licencia}

La opción de CMake \codigo{CMakeLists.txt}{QF_WITH_XAD} (por omisión \texttt{ON}) descarga XAD con \texttt{FetchContent}
y la versión \textbf{2.2.0 fijada por el hash del commit}
(\texttt{aaf4e9cad59cf1bfe9728b491f95f011db687f18}); \codigo{src/CMakeLists.txt}{QF_WITH_XAD} enlaza \texttt{XAD::xad}
con \texttt{qf} (PUBLIC) y define \texttt{QF\_WITH\_XAD}. El bloque activa \texttt{XAD\_USE\_STRONG\_INLINE} (opción de XAD que solo
fuerza \texttt{inline} en sus plantillas de expresiones) y se declara antes de las banderas \texttt{-Wall -Wextra
-Wpedantic} para que no se apliquen al código de XAD. Con la opción en \texttt{OFF} todo lo demás compila y pasa sus pruebas,
\texttt{qfxva.HAS\_XAD} vale \texttt{False} y la ruta responde 501.

XAD se distribuye bajo la \textbf{GNU AGPL-3.0} (o posterior) \citep{xad2026}, y su titular, Xcelerit, ofrece además una licencia
comercial para uso en producción sin las obligaciones de la AGPL. Enlazar una biblioteca AGPL obliga, para quien use el
servicio a través de la red, a ofrecer el código fuente completo de la aplicación salvo que exista la licencia
comercial. El repositorio hoy es público, así que no hay conflicto inmediato; es una decisión de negocio, y la opción
\texttt{QF\_WITH\_XAD} existe para poder prescindir de la biblioteca perdiendo solo estas sensibilidades. La versión
2.2.0 se publicó el 2026-11-11, \emph{el día antes de escribir este capítulo}; se fijó por hash y no por etiqueta porque
la anterior (2.1.0) tiene dos fallos que rompen el patrón de la sección~\ref{sec:cap16-cinta} (\texttt{resetTo}
conservaba las ranuras registradas después de su posición, con lo que el tiempo por trayectoria crecía con el número
de trayectorias, y \texttt{clearDerivativesAfter} descartaba variables registradas antes de su posición), y las corrige 2.2.0 según su
registro de cambios. Al actualizar la biblioteca conviene revisar su estabilidad y repetir las pruebas de la
sección~\ref{sec:cap16-validacion}. Este módulo no entra a la validación contra QuantLib del
capítulo~\ref{cap:validacion}; su validación son las comparaciones con choque y reprecio de esa sección.
\end{sloppypar}

% -------------------------------------------------------------------------------------------------
\section{Ejemplo numérico verificable}\label{sec:cap16-ejemplo}

\begin{sloppypar}
Todos los valores se reproducen con \codigo{docs/tex/verificacion/cap16_xva_aad.py}{main}, que llama al código del
proyecto (\texttt{qfxva} y \texttt{desk.api.risk.xva\_sens}) y, cuando hace falta el desk, usa una base SQLite
temporal con el snapshot MX-DEMO; nunca lee la base real. Necesita \texttt{qfxva} compilado con XAD. Las derivadas
por choque se calculan con \emph{las mismas trayectorias} (misma semilla) que las adjuntas. Las magnitudes marcadas
«recálculo» se rehacen en el script con la biblioteca estándar y el script exige que coincidan con el código. Los
valores de Monte Carlo dependen de la implementación de \texttt{std::normal\_distribution} de \texttt{libstdc++}
(con otra, cambiarían dentro de su error estándar). El script tarda de 30 a 50 segundos, según la carga de la máquina.

\subsection{AAD frente a choque central con las mismas trayectorias}\label{sec:cap16-ej-choque}

Conjunto de compensación de dos swaps: un IRS TIIE-F \emph{pagador} de 100 millones al $6.5\,\%$ (26 periodos de 28
días) y un IRS SOFR \emph{receptor} de 5 millones de \USD\ al $4\,\%$ (8 trimestres de 91 días). Las curvas
son planas, TIIE-F al $7\,\%$ y SOFR al $4\,\%$ (continuo), con pilares en $0.25,0.5,1,2,5,10,30$ años; spot $18.6$;
modelo $a_{\mathrm M}=0.05$, $\sigma_{\mathrm M}=0.01$, $a_{\mathrm U}=0.03$, $\sigma_{\mathrm U}=0.008$,
$\sigma_X=0.12$, $\rho_{\mathrm{MU}}=0.35$, $\rho_{\mathrm{MX}}=0.30$, $\rho_{\mathrm{UX}}=-0.10$; 400 trayectorias, semilla 42 y
malla cada 30 días de 0 a 750 (26 fechas). El crédito de la contraparte es una curva por tramos con vértices en 1 y 3
años e intensidades $1.5\,\%$ y $2.5\,\%$, el propio con $0.6\,\%$ y $1.0\,\%$; $\mathrm{LGD}_c=0.6$,
$\mathrm{LGD}_o=0.5$ y $s_f=0.5\,\%$. La AAD da un CVA de
$\verif{cap16-aad-valor-cva}{11286.2076}$~\MXN\ (tolerancia $5\times10^{-5}$), igual al que da
\texttt{simulate} más \texttt{compute\_xva} con diferencia relativa menor que $10^{-12}$ (el script lo comprueba), y
las mismas derivadas con 1 o con 2 hilos.

La tabla~\ref{tab:cap16-choque} compara, para cinco insumos de clases distintas, la derivada adjunta del CVA con la
diferencia central de \texttt{simulate} y \texttt{compute\_xva}: $\pm0.1$~\pb\ de tasa cero, $\pm10^{-3}$ en el spot,
$\pm10^{-5}$ en $\sigma_{\mathrm M}$ y $\pm10^{-6}$ en la primera intensidad. Los valores llevan 3 decimales (tolerancia absoluta $5\times10^{-4}$).
\begin{table}[ht]
\centering\small
\begin{tabular}{@{}>{\raggedright\arraybackslash}p{6.2cm}rr@{}}
\toprule
Derivada del CVA & AAD & Choque central \\
\midrule
Pilar de 2 años de TIIE-F (\MXN\ por \pb) & \verif{cap16-aad-tiie-2a}{206.278} & \verif{cap16-choque-tiie-2a}{206.276} \\
Pilar de 2 años de SOFR (\MXN\ por \pb) & \verif{cap16-aad-sofr-2a}{-221.855} & \verif{cap16-choque-sofr-2a}{-221.852} \\
Spot (\MXN\ por 1~\%) & \verif{cap16-aad-spot}{34.218} & \verif{cap16-choque-spot}{34.218} \\
$\sigma_{\mathrm M}$ (\MXN\ por \pb\ de vol. normal) & \verif{cap16-aad-sigma}{23.075} & \verif{cap16-choque-sigma}{23.073} \\
Intensidad $h^{\mathrm c}_1$ (\MXN\ por \pb) & \verif{cap16-aad-hazard}{42.606} & \verif{cap16-choque-hazard}{42.606} \\
DV01 paralelo de TIIE-F (suma de pilares) & \verif{cap16-aad-dv01}{129.322} & \verif{cap16-choque-dv01}{129.322} \\
\bottomrule
\end{tabular}
\caption{Derivada adjunta frente a diferencia central con las mismas trayectorias (\codigo{docs/tex/verificacion/cap16_xva_aad.py}{aad_contra_choque}). Todas
las parejas de las tres medidas (15 derivadas) cumplen $|\text{AAD}-\text{choque}|\le\max(1\,\%|\text{choque}|,\,10^{-4})$ (la
regla de las pruebas); la mayor diferencia relativa entre las derivadas de módulo al menos $0.01$ es
\verif{cap16-aad-dif-max}{0.0062}~\% (tolerancia $5\times10^{-5}$).}\label{tab:cap16-choque}
\end{table}
Las dos derivadas coinciden en al menos cuatro cifras significativas; lo que queda es el error de truncamiento de la
diferencia central, y ambas son derivadas del mismo estimador. Los signos son los esperados: el pagador de \MXN\ gana valor, y con él exposición, cuando sube TIIE-F, y el receptor de \USD\ lo
pierde cuando sube SOFR.

\subsection{Sin volatilidad: la derivada exacta}\label{sec:cap16-ej-cero}

Con $\sigma_{\mathrm M}=\sigma_{\mathrm U}=\sigma_X=0$ no hay aleatoriedad y la derivada tiene forma cerrada. Para el
swap de un periodo del capítulo~\ref{cap:xva} (pagador de $\Nom=10^{8}$ al $8.5\,\%$, de 0 a 730 días, curva TIIE-F
plana al $9\,\%$, malla de 10 en 10 días; $\mathrm{LGD}_c=\mathrm{LGD}_o=0.6$, intensidades planas de $2\,\%$ y
$1\,\%$, $s_f=0.5\,\%$), $V/\mathcal B=\Nom(1-F)$ es constante hasta el pago, con
$F=(1+\K\tau)\,\DF{\cTIIE}{T}$, $\tau=730/360$ y $T=2$ años (justo el pilar de 2 años). Como $\partial F/\partial z=-TF$,
y la última fecha con el flujo dentro de $V$ es $t_*=720/365$,
\begin{equation}\label{eq:cap16-cero}
\begin{aligned}
  \mathrm{CVA}&=\mathrm{LGD}_c\,\Nom\,(1-F)\,[1-S_c(t_*)], &
  \frac{\partial\,\mathrm{CVA}}{\partial z^{\mathrm M}_{2a}}&=\mathrm{LGD}_c\,\Nom\,T F\,[1-S_c(t_*)],\\
  \frac{\partial\,\mathrm{FVA}}{\partial z^{\mathrm M}_{2a}}&=s_f\,\Nom\,T F\sum_{i}S_c(t_i)S_o(t_i)\,\Delta t_i, &
  \frac{\partial\,\mathrm{CVA}}{\partial h^{\mathrm c}}&=\mathrm{LGD}_c\,\Nom\,(1-F)\,t_*\,e^{-h^{\mathrm c}t_*},
\end{aligned}
\end{equation}
y todas las demás derivadas respecto de pilares son cero (con interpolación lineal, el pilar de 1 año tiene peso cero), igual que las de SOFR, el spot y el DVA ($V>0$ en
toda trayectoria, así que $V^-=0$). El CVA es $\verif{cap16-cero-valor-cva}{48188.6517}$ (recálculo; la AAD coincide con diferencia
relativa menor que $10^{-12}$). Por \pb\ (recálculo con la biblioteca estándar y AAD; tolerancia $5\times10^{-7}$):
\begin{center}\small
\begin{tabular}{@{}lrr@{}}
\toprule
Derivada & Ecuación~\eqref{eq:cap16-cero} & AAD \\
\midrule
$\partial\mathrm{CVA}/\partial z^{\mathrm M}_{2a}$ & \verif{cap16-cero-dcva-analitica}{454.569749} & \verif{cap16-cero-dcva-aad}{454.569749} \\
$\partial\mathrm{FVA}/\partial z^{\mathrm M}_{2a}$ & \verif{cap16-cero-dfva-analitica}{187.483284} & \verif{cap16-cero-dfva-aad}{187.483284} \\
$\partial\mathrm{CVA}/\partial h^{\mathrm c}$ & \verif{cap16-cero-dhaz-analitica}{236.221656} & \verif{cap16-cero-dhaz-aad}{236.221656} \\
\bottomrule
\end{tabular}
\end{center}
y las demás tasas cero de TIIE-F quedan por debajo de $10^{-9}$ veces la del pilar de 2 años (el script lo exige; son
redondeo de la cancelación entre $\mathcal B$ y $P_{\mathrm M}(t,T)$), y las de SOFR, el spot y el DVA son exactamente cero. Las derivadas respecto de las
propias volatilidades en $\sigma=0$ no se documentan ni se prueban.

\subsection{El estimador es insesgado: la vega frente a la forma cerrada}\label{sec:cap16-ej-vega}

La derivada trayectoria a trayectoria respecto de $\sigma_{\mathrm M}$ no se puede comparar con un choque de las
mismas trayectorias (eso solo prueba que la AAD deriva bien el programa), sino con la derivada \emph{exacta} de algo
que no es Monte Carlo. Para el mismo swap de un periodo (con volatilidad $\sigma_{\mathrm M}=0.01$ y el resto del
modelo del capítulo~\ref{cap:xva}) el CVA, el DVA y el FVA tienen forma cerrada (el recálculo del
capítulo~\ref{cap:xva}, sección «Forma cerrada del CVA, DVA y FVA de un swap de un periodo»), y su derivada respecto de
$\sigma_{\mathrm M}$ se obtiene con una diferencia central de esa función analítica ($\pm10^{-5}$). Se compara con la
media de 20 lotes independientes de la derivada adjunta, 3000 trayectorias por lote y semillas $20261110+7919b$, con su
error estándar entre lotes (\MXN\ por \pb\ de volatilidad normal; tolerancia $5\times10^{-5}$):
\begin{center}\small
\begin{tabular}{@{}lrrr@{}}
\toprule
Medida & Exacta & Media AAD & Error estándar \\
\midrule
CVA & \verif{cap16-vega-exacta-cva}{39.3950} & \verif{cap16-vega-media-cva}{37.2346} & \verif{cap16-vega-ee-cva}{0.9575} \\
DVA & \verif{cap16-vega-exacta-dva}{19.9478} & \verif{cap16-vega-media-dva}{19.7411} & \verif{cap16-vega-ee-dva}{0.2320} \\
FVA & \verif{cap16-vega-exacta-fva}{0.0000} & \verif{cap16-vega-media-fva}{-0.7214} & \verif{cap16-vega-ee-fva}{0.5154} \\
\bottomrule
\end{tabular}
\end{center}
El mayor desvío es de $\verif{cap16-vega-z-max}{2.256}$ errores estándar (tolerancia $5\times10^{-4}$; el script exige
menos de 3), el del CVA, que queda del lado bajo. Con estas semillas el CVA queda por debajo de la derivada exacta de forma persistente, pero
los tamaños de lote que se probaron comparten trayectorias y no son evidencia independiente. Dos corridas aparte, con otras
semillas y la misma malla, no incluidas en el script por su duración (de 1 a 2 minutos), no muestran sesgo: 40 lotes de 5000 trayectorias (semillas $777+31b$) dieron para el CVA una
media de $39.97\pm0.56$ frente a $39.395$ exacta ($z=1.03$) y para el DVA $20.11\pm0.13$ frente a $19.948$ ($z=1.25$); 60
lotes de 3000 (semillas $424242+104729b$), $38.98\pm0.46$ ($z=-0.89$) y $19.96\pm0.15$ ($z=0.07$). El FVA no depende de
$\sigma_{\mathrm M}$ en valor esperado ($\E[V/\mathcal B]$ es constante: martingala) y su estimador, de media cercana a
cero, lo reproduce ($-0.18\pm0.26$ en la segunda corrida). Del error estándar entre lotes se obtiene lo que cuesta el
ruido: la desviación entre corridas de 3000 trayectorias es $0.9575\sqrt{20}\approx4.28$ (un $10.9\,\%$ del valor exacto)
para la vega del CVA y $0.2320\sqrt{20}\approx1.04$ ($5.2\,\%$) para la del DVA, y crece como $\sqrt{n}$ al reducir
las trayectorias (unos $13\,\%$ con las 2000 de la configuración por omisión). La vega por trayectoria es mucho más
ruidosa que el CVA mismo, cuyo error estándar para este swap es de $0.37\,\%$ con 20\,000 trayectorias en el
capítulo~\ref{cap:xva}, es decir, cerca de $1\,\%$ con 3000; y la ruta no informa el error estándar de las derivadas.

\subsection{Crédito: del jacobiano a las cotizaciones de CDS, y la ruta}\label{sec:cap16-ej-credito}

Banco Ejemplo tiene un IRS TIIE-F pagador de 200 millones al $6.5\,\%$ de 26 periodos, pactado el 2026-11-10. Su curva de crédito
sintética es de $90,100,110,125,135,150$~\pb\ a 1, 2, 3, 5, 7 y 10 años con $R=40\,\%$, y la propia de
$70,\dots,125$~\pb; $s_f=0.5\,\%$, 400 trayectorias y semilla 42. La ruta (\texttt{xva\_sens.compute}) devuelve un CVA de
referencia de $\verif{cap16-api-cva}{12281.9218}$~\MXN, igual a igual al de \texttt{GET /api/xva} (el script
compara con \texttt{==} las cuatro cifras), con $\verif{cap16-api-n-insumos}{48}$ insumos. Las derivadas del CVA respecto de la
intensidad de los dos primeros tramos de la contraparte (\MXN\ por \pb\ de intensidad) son
$\verif{cap16-cred-dcva-h1}{48.4445}$ y $\verif{cap16-cred-dcva-h2}{25.9164}$; el resto es exactamente cero porque el
swap vence antes de 2 años y la exposición no llega a los tramos de 3 años o más. El jacobiano del bootstrap por
reprecio, ecuación~\eqref{eq:cap16-jacobiano-reprecio}, es triangular (el script comprueba que $J_{ki}=0$ para $i<k$) y
tiene
\[
  J_{11}=\verif{cap16-cred-j11}{1.675581},\qquad
  J_{12}=\verif{cap16-cred-j12}{-1.822996},\qquad
  J_{22}=\verif{cap16-cred-j22}{3.501813}
\]
(tolerancia $5\times10^{-7}$; $J_{11}$ es cercano a la «regla del triángulo» $1/(1-R)=1.667$). Un punto base más en la
cotización de 1 año sube la intensidad del primer tramo en $1.68$~\pb\ pero \emph{baja} la del segundo en $1.82$~\pb: el
CDS de 2 años, cuya cotización no se movió, debe seguir valiendo cero con la supervivencia más baja del primer año. La regla
de la cadena, ecuación~\eqref{eq:cap16-cadena}, da entonces para el CDS de 1 año $48.4445\cdot1.675581+25.9164\cdot(-1.822996)$ y,
para el de 2 años, $25.9164\cdot3.501813$ (el script exige que coincida con la ruta a $10^{-9}$ relativo); las cotizaciones de 3 años en adelante
valen exactamente cero. Contra el reprecio del CDS (la cotización movida $\pm0.01$~\pb, el bootstrap vuelto a correr y el XVA
recalculado con el mismo perfil), en \MXN\ por \pb\ de cotización (tolerancia $5\times10^{-7}$):
\begin{center}\small
\begin{tabular}{@{}lrr@{}}
\toprule
Derivada del CVA & Ruta (cadena) & Reprecio del CDS \\
\midrule
Cotización de 1 año & \verif{cap16-cred-dcva-cds1-api}{33.927133} & \verif{cap16-cred-dcva-cds1-choque}{33.927133} \\
Cotización de 2 años & \verif{cap16-cred-dcva-cds2-api}{90.754376} & \verif{cap16-cred-dcva-cds2-choque}{90.754376} \\
\bottomrule
\end{tabular}
\end{center}
Como comprobación independiente de la cinta, el script rehace la ecuación~\eqref{eq:cap16-dcva-dh} con el perfil
$\mathrm{EPE}^*$ de \texttt{simulate} y la supervivencia por tramos en la biblioteca estándar, y coincide con los
$\partial\mathrm{CVA}/\partial h^{\mathrm c}_i$ de la ruta a $10^{-9}$ relativo. Finalmente, el FVA es lineal en el diferencial de
fondeo, $\partial\mathrm{FVA}/\partial s_f=\mathrm{FVA}/s_f$ (ecuación~\eqref{eq:cap16-dcva-dh}): la ruta da
$\verif{cap16-cred-fondeo}{28.1507}$~\MXN\ por \pb\ de fondeo, igual a $10^{-4}\,\mathrm{FVA}/s_f$ (a $10^{-9}$ relativo).

\subsection{Derivada parcial frente a derivada total}\label{sec:cap16-ej-parcial}

Banco Cruzado tiene un cross currency swap que paga \USD\ por 10 millones de \USD\ a 18.2, con diferencial de $5$~\pb,
26 periodos e intercambios inicial y final de nocional, pactado el 2026-11-10; su pata \MXN\ se descuenta con MXN-USD. El
snapshot MX-DEMO calibra MXN-USD con puntos forward (pilares de 1 mes a 1 año) y con cotizaciones de base sobre TIIE-F
(pilares de 2 a 10 años), de modo que un desplazamiento de TIIE-F la mueve. El desk sabe hacerlo
(\codigo{desk/api/market.py}{MarketContext.shifted} recalibra MXN-USD con sus cotizaciones fijas), y eso permite medir las dos derivadas
del CVA respecto de un desplazamiento paralelo de $+1$~\pb\ de TIIE-F, con las mismas trayectorias, ecuación~\eqref{eq:cap16-total}
(\MXN\ por \pb; tolerancia $5\times10^{-4}$):
\begin{center}\small
\begin{tabular}{@{}>{\raggedright\arraybackslash}p{7.0cm}rr@{}}
\toprule
DV01 de TIIE-F & CVA & Total ($\mathrm{CVA}-\mathrm{DVA}+\mathrm{FVA}$) \\
\midrule
Parcial por AAD (suma de los 10 pilares) & \verif{cap16-parcial-aad-cva}{111.858} & \verif{cap16-parcial-aad-total}{470.996} \\
Parcial por choque (MXN-USD fija) & \verif{cap16-parcial-choque-cva}{111.855} & \verif{cap16-parcial-choque-total}{470.995} \\
Total por choque (MXN-USD recalibrada) & \verif{cap16-total-choque-cva}{-145.508} & \verif{cap16-total-choque-total}{-458.339} \\
Total por la API (\texttt{recalibrated}, sin simular más) & \verif{cap16-total-api-cva}{-145.507} & \verif{cap16-total-api-total}{-458.339} \\
\bottomrule
\end{tabular}
\end{center}
La AAD por sí sola reproduce la derivada \emph{parcial}; la total tiene el signo contrario. La diferencia es exactamente el segundo
sumando de la ecuación~\eqref{eq:cap16-total}: los pilares de MXN-USD de 2 años en adelante se mueven $\verif{cap16-mueve-2a}{1.0400}$~\pb\ (el de 2 años) a
$\verif{cap16-mueve-10a}{1.0054}$~\pb\ (el de 10 años) por cada punto base de TIIE-F, y los cuatro primeros no se mueven; sumando
a la parcial los productos de las sensibilidades \emph{por AAD} a los pilares de MXN-USD por esos movimientos se
reconstruye la total del CVA, $\verif{cap16-total-reconstruida-cva}{-145.507}$ (a menos de $10^{-5}$ relativo de la medida por choque y recalibración; el script exige $1\,\%$). Eso es
lo que hace la ruta, con el jacobiano por pilar en lugar del del desplazamiento paralelo: su fila \texttt{recalibrated}
(última de la tabla) coincide con el choque con recalibración a menos del $1\,\%$ que exige el script, y con la
reconstrucción a menos de $10^{-3}$ relativo. El crédito se mantuvo fijo en los
cuatro renglones (no se recalibró el bootstrap de CDS, que también descuenta con TIIE-F). Para una cartera cuyas
operaciones se descuentan con MXN-USD, el DV01 \emph{parcial} de TIIE-F \textbf{no es el cambio del XVA ante un movimiento
de 1~\pb\ de TIIE-F en el mercado}: ese es el total, y es el que sirve para cubrir.
\end{sloppypar}

% -------------------------------------------------------------------------------------------------
\section{Validación}\label{sec:cap16-validacion}

\begin{sloppypar}
\subsection{Pruebas del repositorio}\label{sec:cap16-pruebas}

\paragraph{Corrección de las derivadas (C++).} La regla de aceptación del corte es que cada derivada adjunta
coincida con la diferencia central con las mismas trayectorias dentro de $\max(1\,\%,\,1\text{ \MXN\ por unidad del insumo})$
(que es $10^{-4}$ \MXN\ por \pb), sin aflojar. Cuatro pruebas cubren todas las clases de insumo con 400 trayectorias:
\codigo{tests/test_ExposureEngine.cpp}{AadMatchesBumpTiieSwap} (un swap TIIE-F; además exige que $\partial/\partial$SOFR,
MXN-USD y spot sean exactamente cero),
\codigo{tests/test_ExposureEngine.cpp}{AadMatchesBumpSofrSwap} (un swap SOFR, que depende de las tres curvas y del spot),
\codigo{tests/test_ExposureEngine.cpp}{AadMatchesBumpXccySwap} y
\codigo{tests/test_ExposureEngine.cpp}{AadMatchesBumpCombinedSet} (el conjunto de cuatro swaps y tres cross currency
swaps, con las tres curvas, el spot, los cinco números del modelo, las intensidades de ambas curvas y el fondeo). El
desacuerdo medido es de $\sim10^{-4}$ relativo en curvas, modelo y spot (el máximo, $4.5\times10^{-3}$, en un pilar de
derivada pequeña: $-\num{22542}$ frente a $-\num{22644}$) y de $10^{-10}$ en crédito y fondeo.
\codigo{tests/test_ExposureEngine.cpp}{AadValueMatchesComputeXva} exige que el valor sea el de \texttt{computeXva}
(relativo $10^{-12}$, y los errores estándar);
\codigo{tests/test_ExposureEngine.cpp}{AadZeroVolatilityMatchesDeterministicDerivative} compara con la derivada determinista
de un choque de $10^{-7}$ (tolerancia $10^{-6}$ relativa más $10^{-3}$ de redondeo del choque);
\codigo{tests/test_ExposureEngine.cpp}{AadFlatHazardHasOneTranche},
\codigo{tests/test_ExposureEngine.cpp}{AadRejectsNonLinearCurves} (interpolación no lineal y tipos de curva de crédito
no admitidos) y \codigo{tests/test_ExposureEngine.cpp}{AadMaxThreadsCapsTheWorkersAndKeepsTheResult} (el tope de hilos
acota las cintas y el resultado cambia solo por redondeo, $10^{-9}$). Sin XAD,
\codigo{tests/test_ExposureEngine.cpp}{XvaSensitivitiesWithoutXadThrowsLogicError} comprueba el \texttt{logic\_error}
(se corrió en una compilación aparte con la opción en \texttt{OFF}: todas las pruebas pasan). Bindings:
\codigo{tests/test_xva_engine_py.py}{test_has_xad_flag_is_a_bool},
\codigo{tests/test_xva_engine_py.py}{test_xva_sensitivities_shape} y
\codigo{tests/test_xva_engine_py.py}{test_xva_sensitivities_max_threads}.

\paragraph{Por la API.} \codigo{desk/api/tests/test_xva_sens.py}{test_aad_matches_bump_and_reprice_for_one_input_of_each_class}
repite la comparación sobre el conjunto sembrado de una contraparte, a través de la ruta, para un insumo de cada
clase: el pilar de mayor sensibilidad de TIIE-F, SOFR y MXN-USD, el spot, la volatilidad de \USD/\MXN, el fondeo y
una cotización de CDS (con el bootstrap vuelto a correr); \codigo{desk/api/tests/test_xva_sens.py}{test_xva_is_exactly_the_one_of_get_xva_and_dv01_sums_the_pillars}
exige que el XVA de referencia sea el de \texttt{GET /api/xva} y que cada DV01 (parcial y total) sea la suma de sus pilares;
\codigo{desk/api/tests/test_xva_sens.py}{test_total_matches_bump_recalibrate_and_reprice} compara la total con el
choque del insumo, la recalibración de MXN-USD con \texttt{MarketContext.shifted} y el reprecio completo con las mismas
500 trayectorias, para un pilar de TIIE-F, uno de SOFR y el spot (dentro de $\max(1\,\%,\,1~\MXN\text{ por }\pb)$), y exige
además que lo que la recalibración aporta (total menos parcial) coincida con el de ese reprecio dentro del $2\,\%$, lo que
delata un jacobiano mal escalado (tres mutaciones hechas a mano y revertidas, el ancho del jacobiano de TIIE-F, la escala
del spot y la omisión de la recalibración de TIIE-F, hacen fallar la prueba); los
tipos, longitudes y totales, en \codigo{desk/api/tests/test_xva_sens.py}{test_shape_lengths_and_totals}; los
errores, en \codigo{desk/api/tests/test_xva_sens.py}{test_unknown_counterparty_is_404_and_missing_parameters_are_400},
\codigo{desk/api/tests/test_xva_sens.py}{test_without_xad_the_route_answers_501},
\codigo{desk/api/tests/test_xva_sens.py}{test_counterparty_without_covered_trades_has_empty_sensitivities} y
\codigo{desk/api/tests/test_xva_sens.py}{test_without_the_mxn_usd_curve_it_does_not_appear}.

\paragraph{Pruebas por mutación (a mano, no automáticas).} Para comprobar que las pruebas detectan un adjunto mal
propagado se sustituyó, en la interpolación de pilares, \texttt{z[i+1]} por \texttt{Real(xad::value(z[i+1]))}
(se pierde la dependencia del pilar superior de cada intervalo). Fallaron las pruebas del swap TIIE-F, del SOFR, del
cross currency swap, del conjunto combinado y del límite sin volatilidad (por ejemplo, $\partial\mathrm{CVA}/\partial z_{\mathrm{TIIE}}(1a)$:
AAD $-\num{1049708}$ frente a $-\num{1379976}$ del choque), mientras la prueba dorada en \texttt{double} siguió pasando: la
mutación solo afecta al tipo activo. Para la ruta se comprobaron además tres mutaciones (omitir el jacobiano del CDS,
duplicar la escala del spot y desalinear un pilar), que la comparación con el reprecio detectó. La mutación fue
revertida.

\subsection{Costo y memoria medidos}\label{sec:cap16-costo}

Las cifras de esta sección son \textbf{mediciones de una sola máquina} (un portátil i7-10510U de 4 núcleos y 8 hilos,
con carga ajena de 1 a 7 tareas durante las corridas) \textbf{y de la compilación de depuración} (\texttt{CMAKE\_BUILD\_TYPE=Debug}
de \texttt{build/}), que es la que sirve la API: con una compilación optimizada los tiempos absolutos y las razones serían
otros, y no se midieron. Son órdenes de magnitud, con unos $\pm20\,\%$; no se verifican en el script (varían). Se citan las cifras medidas durante el desarrollo con la prueba desactivada
\codigo{tests/test_ExposureEngine.cpp}{DISABLED_AadTiming} (se corre a mano), con un conjunto pesado: tres operaciones
(IRS TIIE-F pagador a 10 años con 130 periodos, IRS SOFR a 5 años, cross currency swap a 5 años), malla de integración de
1407 fechas y 2000 trayectorias.

\begin{table}[ht]
\centering\small
\begin{tabular}{@{}lrrr@{}}
\toprule
Hilos & Valuación \texttt{double} (s) & AAD (s) & Razón \\
\midrule
1 & 4.87 & 44.5 & 9.1 \\
2 & 2.83 & 27.7 & 9.8 \\
4 & 1.54 & 26.0 & 16.8 \\
8 (valor por omisión de \texttt{simulate}) & 1.31 & 26.4 & 20.2 \\
\midrule
1 núcleo fijado para ambas (afinidad), 2000 trayectorias & 4.80 & 44.5 & 9.27 \\
1 núcleo fijado para ambas, 500 trayectorias & 1.11 & 10.9 & 9.85 \\
Como la API: valuación en los 8 hilos, AAD en 2 & 1.28 & 27.6 & 21.6 \\
\bottomrule
\end{tabular}
\caption{Costo de las sensibilidades por AAD frente a una valuación (\texttt{simulate} más \texttt{computeXva}) del mismo
conjunto, con la misma malla; filas superiores, el mismo número de hilos en las dos columnas; las dos últimas, la
razón con un solo núcleo para ambas (el criterio de aceptación del corte) y la que ve el usuario de la API.}\label{tab:cap16-costo}
\end{table}

\paragraph{A igual número de hilos la razón es de 7 a 9 valuaciones.} Con un núcleo para ambas, la AAD cuesta
\textbf{9.27 valuaciones} con 2000 trayectorias (9.1 en la primera medición, 9.85 con 500 trayectorias: el margen
frente al objetivo de 10 es corto) en el conjunto pesado; una medición independiente sobre las tres contrapartes
sembradas, con un núcleo para ambas (el mejor de tres corridas), dio 7.4, 8.0 y 6.5 en Banco Alfa, Beta y Gamma. De 6.5 a
9.9, es decir, \textbf{de 7 a 9 valuaciones} según el conjunto. Con los 48 insumos del snapshot, los choques serían $2\times48=96$ valuaciones
($\approx10$ veces más). La razón \textbf{no se mantiene al repartir en hilos}: la valuación en \texttt{double} escala
(3.7 veces con 8 hilos) y la AAD no pasa de 1.7 veces (de unos 44 s a unos 26 s): es una carga limitada por el ancho de banda
de memoria, no por cómputo. El barrido de adjuntos de un hilo pasa de 17.8 a 115 s de CPU con 8 hilos (6.5 veces) y la
grabación de 24.1 a 82 s (3.4 veces), mientras la valuación en \texttt{double} casi no se infla (9.4 contra 9.5 s). Por
eso la API usa como máximo 2 hilos (más no acelera y multiplica la memoria), y con ellos contra la valuación en los 8
hilos del equipo la razón es de unas $21$ veces (última fila). Comparada con la estimación por choques
(96 valuaciones en los 8 hilos, 125 s para ese conjunto, estimada y no medida), la AAD de 26 a 28 s sigue siendo unas 5 veces
más barata.

\paragraph{Dónde se va el tiempo} (un hilo, 2000 trayectorias, tiempos de CPU). Valuación en \texttt{double}: 4.8 s
(trayectoria 1.15, bonos 2.45, \texttt{setValue} y resto 1.17). AAD: 24.1 s de grabación (trayectoria 3.2, \texttt{computeBonds}
7.6, \texttt{setValue} 13.1, resto 0.3) y 17.8 s (el 40\,\%) de barrido de adjuntos. Los bonos y la valuación de patas
(la fórmula afín) son el 70\,\% de la cinta. Un adjunto escrito a mano de esa fórmula evitaría grabar unos 2.4 millones de
operaciones por trayectoria y es la mejora futura señalada (sección~\ref{sec:cap16-limites}).

\paragraph{Memoria.} El pico de una cinta por hilo es de $\num{48063764}$ bytes (48.06 MB) \emph{y es idéntico con 500 y con
2000 trayectorias}, como dice la ecuación~\eqref{eq:cap16-memoria}. La parte común ocupa $\num{4957880}$ bytes (4.96 MB) y
cada trayectoria agrega unos $\num{43105884}$ bytes (43.1 MB) que \texttt{resetTo} descarta: una cinta de toda la simulación,
con 2000 trayectorias, sería de unos 86 GB. Con los dos hilos de la API, unos 96 MB de cinta. (Los 5--6 MB de la sonda
original medían solo la parte común, después del reinicio.)

\paragraph{Tiempo real por la API.} La tabla~\ref{tab:cap16-api} da los tiempos de las tres contrapartes sembradas de
MX-DEMO (2000 trayectorias, 48 insumos). \texttt{valuation\_elapsed\_s} incluye armar el conjunto, el bootstrap de crédito y el
perfil publicado, de modo que sobreestima un poco la valuación pura, y la estimación por choques
(\texttt{bump\_cost\_estimate\_s}) supone los choques con esa valuación, es decir, en todos los hilos.
\begin{table}[ht]
\centering\small
\begin{tabular}{@{}lrrrrr@{}}
\toprule
Contraparte & AAD (s) & Valuación (s) & Razón & Choques estimados (s) & Total ruta (s) \\
\midrule
\multicolumn{6}{@{}l}{\emph{Medición del desarrollo} (carga ajena $\approx1.6$)} \\
Banco Alfa & 5.77 & 0.31 & 18.6 & 30.0 & 6.1 \\
Banco Beta & 11.77 & 0.60 & 19.6 & 57.5 & 12.4 \\
Banco Gamma & 5.27 & 0.35 & 15.1 & 33.3 & 5.6 \\
\midrule
\multicolumn{6}{@{}l}{\emph{Repetición al escribir este capítulo} (carga ajena $\approx1.2$)} \\
Banco Alfa & 8.23 & 0.38 & 21.8 & 36.2 & 8.6 \\
Banco Beta & 14.25 & 0.61 & 23.2 & 59.0 & 14.9 \\
Banco Gamma & 9.43 & 0.52 & 18.1 & 50.0 & 10.0 \\
\bottomrule
\end{tabular}
\caption{Tiempos reales de \texttt{GET /api/xva/sensitivities} para las tres contrapartes sembradas (2000 trayectorias). La razón es
\texttt{elapsed\_s}$/$\texttt{valuation\_elapsed\_s}; los choques son $2\times48\times$\texttt{valuation\_elapsed\_s}, una
estimación (no se corrieron).}\label{tab:cap16-api}
\end{table}
Con la malla fina y 2000 trayectorias los conjuntos sembrados tardan de unos 5 a 15 s, no los 26 s del conjunto pesado,
y la AAD de las tres es de 5 a 6 veces más barata que la estimación por choques. Repetido al añadir la derivada total
(que suma unos 0.07~s de recalibraciones al \texttt{elapsed\_s}): Alfa 5.76~s contra 0.31 (razón 18.7), Beta 11.81 contra
0.61 (19.4) y Gamma 5.23 contra 0.34 (15.5), con choques estimados de 29.5, 58.3 y 32.5~s. La razón que ve el usuario (15 a 23) es mayor que
la de 9 a 10 de la tabla~\ref{tab:cap16-costo} porque la valuación corre en todos los hilos y la AAD en dos (hallazgo H25).
\end{sloppypar}

% -------------------------------------------------------------------------------------------------
\section{Límites y supuestos}\label{sec:cap16-limites}

\begin{limites}
  \begin{itemize}
    \item \textbf{Parcial y total.} La AAD entrega la derivada parcial (MXN-USD y crédito fijos). La ruta añade la total
          con MXN-USD recalibrada, con sus puntos forward y su base fijos (que no se derivan), para los pilares de TIIE-F y de
          SOFR y el spot, y es la que sirve para cubrir: puede tener el signo contrario al de la parcial
          (sección~\ref{sec:cap16-ej-parcial}). \emph{Ninguna} recalibra las curvas de crédito (el bootstrap de CDS
          descuenta con TIIE-F): faltan los dos últimos sumandos de la ecuación~\eqref{eq:cap16-total}. MXN-USD no tiene
          total (es lo que se recalibra), y la de las volatilidades, el crédito y el fondeo es la parcial. Las curvas de tasa se derivan por tasa cero del pilar, no por
          cotización de depósito u OIS; solo el crédito pasa a cotizaciones (por el jacobiano del bootstrap
          obtenido por reprecio, no por AAD del bootstrap).
    \item \textbf{Primer orden, sin correlaciones y sin PFE.} No hay segundo orden (gamma cruzada), ni sensibilidad a
          las correlaciones, a la recuperación (LGD) ni a las fechas y términos de las operaciones; no se derivan el
          PFE (un cuantil) ni el XVA incremental. Se deriva una contraparte (un conjunto de compensación) a la vez.
    \item \textbf{Derivadas de la estimación, con su ruido.} Son derivadas trayectoria a trayectoria con las mismas
          trayectorias, malla de integración y semilla que el XVA; llevan el error de Monte Carlo del XVA y, la vega,
          bastante más (11\,\% entre corridas de 3000 trayectorias en el ejemplo de la
          sección~\ref{sec:cap16-ej-vega}, contra cerca de $1\,\%$ del CVA). La ruta \emph{no informa el error estándar de las derivadas}. Son las de la malla $\mathcal G$ con
          extremo derecho (capítulo~\ref{cap:xva}), no las del XVA en tiempo continuo.
    \item \textbf{Tiempo.} A igual número de hilos, de 7 a 9 valuaciones; en la ruta, con la valuación en todos los
          hilos y la adjunta en dos, de 15 a 23 veces el tiempo de una valuación, y decenas de segundos en conjuntos
          largos (26 s el conjunto de 3 operaciones a 10 años). Es 5 veces más barata que los choques ($2m$
          valuaciones, que se estiman con la valuación en todos los hilos y no se corrieron) con 48 insumos, pero no
          instantánea; por eso se calcula a petición. Todas las cifras son de la compilación de depuración que sirve
          la API.
    \item \textbf{Memoria.} Unos 48 MB de cinta por hilo con la malla fina (4.96 MB de parte común y 43 MB por la
          trayectoria en curso), independientes del número de trayectorias, pero proporcionales al tamaño del conjunto y
          a la malla; con más hilos, proporcionalmente más. El tope de 2 hilos acota la memoria.
    \item \textbf{Curvas admitidas.} Solo curvas de tasa con interpolación \texttt{Linear} y solo curvas de crédito
          \texttt{PiecewiseHazardCurve} o \texttt{FlatHazardRate} (otras lanzan \texttt{std::invalid\_argument}).
    \item \textbf{Volatilidad exactamente cero.} Con $\sigma_X$ o $\sigma_c$ exactamente en cero su sensibilidad sale
          \emph{cero}, porque la factorización de Cholesky semidefinida
          (\codigo{src/xva/exposure_engine.cpp}{choleskySemidefinite}) anula la columna correspondiente: no es la derivada
          lateral. No afecta a los snapshots reales, cuyas volatilidades son positivas.
    \item \textbf{Licencia.} XAD es AGPL-3.0 (con licencia comercial aparte). Un producto que se ofrezca a un cliente sin
          publicar su código necesita esa licencia o compilar con \texttt{QF\_WITH\_XAD=OFF}, que deja la ruta en 501.
          Es una decisión de negocio pendiente.
    \item \textbf{Versión de XAD.} 2.2.0 se publicó el 2026-11-11, un día antes de este capítulo, y está fijada por hash
          de commit. Se aceptó porque corrige justo la técnica que se usa y las pruebas de corrección contra choques
          validan los números; conviene revisar su estabilidad al actualizar.
    \item \textbf{Mejora futura.} Escribir a mano el adjunto de la fórmula afín de los bonos y las patas (el 70\,\% de
          la cinta) reduciría la cinta y el tiempo; con un adjunto de tamaño 1 (solo el total) la cinta es un 20\,\% más
          rápida. Ninguna está hecha. Las cifras de costo son de una máquina y varían con la carga.
    \item \textbf{Informativo.} Como el XVA, las sensibilidades no entran al valor presente, a los escenarios ni a la
          contabilidad de coberturas, y no están contrastadas con QuantLib.
  \end{itemize}
\end{limites}

% -------------------------------------------------------------------------------------------------
\section{Referencias}\label{sec:cap16-referencias}

El modo adjunto y su costo, a \citet{griewank2008}, con una introducción práctica en \citet{naumann2011}; el uso del
adjunto para sensibilidades por Monte Carlo en finanzas, a \citet{giles2005} (que presenta el método
trayectoria a trayectoria recorrido hacia atrás) y \citet{capriotti2011}, su aplicación al riesgo de crédito de
contraparte, a \citet{capriotti2011b}, y una revisión reciente, a \citet{capriotti2024}; el patrón de cinta por
trayectoria y por hilo, a \citet{savine2018}; el estimador de la derivada por trayectoria y su insesgadez, a
\citet{broadie1996} y \citet[cap.~7]{glasserman2003}; y la biblioteca, su versión y su licencia, a \citet{xad2026}.
El motor que se deriva y sus convenciones son los del capítulo~\ref{cap:xva}.
